\documentclass[11pt,a4paper]{amsart}
\pdfinfoomitdate=1
\pdftrailerid{}
\pdfsuppressptexinfo=15
\usepackage[T1]{fontenc}
\usepackage{lmodern}
\usepackage[margin=27mm]{geometry}
\usepackage{microtype}
\usepackage{amsmath,amssymb,mathtools}
\usepackage{booktabs}
\usepackage{array}
\usepackage{xcolor}
\usepackage[numbers,sort&compress]{natbib}
\usepackage[colorlinks=true,linkcolor=blue!55!black,citecolor=blue!55!black,urlcolor=blue!55!black]{hyperref}
\usepackage[nameinlink,noabbrev]{cleveref}
\raggedbottom

\newtheorem{theorem}{Theorem}[section]
\newtheorem{proposition}[theorem]{Proposition}
\newtheorem{corollary}[theorem]{Corollary}
\newtheorem{lemma}[theorem]{Lemma}
\theoremstyle{definition}
\newtheorem{definition}[theorem]{Definition}
\theoremstyle{remark}
\newtheorem{remark}[theorem]{Remark}

\newcommand{\F}{\mathbb F}
\newcommand{\Z}{\mathbb Z}
\newcommand{\N}{\mathbb N}
\newcommand{\Prob}{\mathbb P}
\newcommand{\E}{\mathbb E}
\newcommand{\cL}{\mathcal L}
\newcommand{\cX}{\mathcal X}
\newcommand{\supp}{\operatorname{supp}}
\newcommand{\htop}{h_{\mathrm{top}}}
\newcommand{\hb}{h_{\mathrm b}}
\newcommand{\one}{\mathbf 1}

\title[Finite-field adjacent-product processes]{Finite-Field Adjacent-Product Processes:\\
One-Dependence, Infinite Markov Memory, and Exact Support Entropy}
\author{Anonymous}
\date{Internal Stage 2 draft, 28 August 2026}
\subjclass[2020]{37A35, 37B10, 60G10, 60J10, 94A17}
\keywords{finite field, hidden Markov process, one-dependent process, variable-length memory, shift of finite type, entropy rate}

\begin{document}

\begin{abstract}
Let $(U_i)_{i\in\Z}$ be independent and uniform on the finite field
$\F_q$, and observe only the adjacent products
$Y_i=U_iU_{i+1}$.  This elementary two-block factor exhibits a sharp
separation between dependence range and prediction memory.  Its law is
stationary, reversible, mixing, and one-dependent, but it is not a Markov
chain of any finite order.  We identify its support as the mixing two-step
shift of finite type obtained by forbidding $a0b$ with $a,b\ne0$.  If $L_n$
is the number of supported words of length $n$, then, with $m=q-1$,
\[
 L_{n+3}=qL_{n+2}-mL_{n+1}+mL_n,
 \qquad
 (L_0,L_1,L_2)=(1,q,q^2).
\]
Consequently the support entropy is $\log\lambda_q$, where $\lambda_q$ is
the Perron root of
$t^3-qt^2+(q-1)t-(q-1)$.  Exact $2\times2$ matrices give every cylinder
probability.  More sharply, if a nonzero observation is followed by $r$
zeros, then, for $r\geq1$ and every fixed $b\ne0$, the next-symbol probability is
\[
 \Prob(Y_0=b\mid a0^r)
 =\frac{(q-1)F_{r-1}}{qF_{r+1}},
 \qquad F_{r+2}=F_{r+1}+(q-1)F_r.
\]
These probabilities oscillate forever and prove infinite Markov order.
The same renewal contexts yield an explicit exponentially convergent series
for the measure entropy.  We state the hidden-Markov and variable-memory
ownership boundary and make no absolute priority claim.
\end{abstract}

\maketitle

\section{Introduction}

A finite-range factor of an independent process need not have finite
prediction memory.  The distinction is easy to miss: one-dependence says
that sufficiently separated \emph{sets} of coordinates are independent,
whereas a finite-order Markov property asks whether the most recent bounded
window screens the future from the entire past.  The process studied here
separates the two notions within a minimal algebraic model.

Fix a prime power $q\geq2$.  From independent uniform vertices
$U_i\in\F_q$, put on every edge of the integer line the product
$Y_i=U_iU_{i+1}$.  Nonzero products propagate invertibility along a run,
while a zero records only that at least one endpoint vanished.  The field
zero is therefore both a local singularity and an information-losing reset.
That single rank drop is responsible for all of the structure below.

Two-block factors of iid sources are elementary examples of one-dependent
processes, inside a literature that includes Aaronson, Gilat, Keane, and de
Valk's study of the converse block-factor problem \cite{AaronsonEtAl1989}.
We take that general relationship as prior.  Functions of finite-state Markov
chains and their entropy go back at least to Blackwell \cite{Blackwell1957}.
Context models with variable memory were
introduced in information theory by Rissanen \cite{Rissanen1983} and
formalized statistically by B\"uhlmann and Wyner
\cite{BuhlmannWyner1999}.  Our process is a finite-state hidden Markov process,
and its prediction law has an almost-surely finite but unbounded context-tree
description; we use the variable-memory comparison in this broad sense, not
as a claim of membership in a finite-context statistical subclass.  The
contribution of this note is the exact finite-field specialization:
the support, all cylinder fibers, the unbounded context law, and the entropy
series are obtained in closed form from the same second-order recurrence.

The main results are as follows.
\begin{enumerate}
\item The topological support is exactly the two-step shift forbidding the
      words $a0b$ with $a,b\in\F_q^\times$.  A four-state transfer matrix
      gives the exact word recurrence and its Perron entropy.
\item The measure is reversible and one-dependent, hence strongly mixing,
      but has no finite Markov order.  Its prediction after a zero run is a
      generalized-Fibonacci ratio that alternates around a nontrivial
      limit and never stabilizes.
\item Two $2\times2$ matrices count every finite fiber.  They also give an
      explicit context expansion for the entropy rate and show that the
      measure has a strict entropy gap below the support entropy.
\end{enumerate}

All statements hold for every finite field, including nonprime fields.  The
proofs use only that multiplication by a nonzero field element is a
bijection.  The accompanying program checks the defining map, the forbidden
language, the fiber matrices, the recurrence, the prediction ratios, and the
entropy series over $\F_2$, $\F_3$, $\F_4$, and $\F_5$.

\section{The process and its finite fibers}
\label{sec:setup}

Let $q$ be a prime power, let $\F_q$ be the field with $q$ elements, and set
$m=q-1$.  On the Bernoulli system
$(\F_q^\Z,\operatorname{Unif}(\F_q)^\Z,\sigma)$ define the two-block code
\begin{equation}\label{eq:factor}
 \Phi(u)_i=u_i u_{i+1},\qquad i\in\Z.
\end{equation}
Write $Y=\Phi(U)$, let $\mu_q$ be its law, and let
$X_q=\supp(\mu_q)\subseteq\F_q^\Z$.

We first record a fiber formula that will be used in each later section.  A
partial hidden word ending at a vertex can end either at zero or at a
nonzero field element.  In that order define
\begin{equation}\label{eq:fiber-matrices}
 C_0=\begin{pmatrix}1&m\\[2pt]1&0\end{pmatrix},
 \qquad
 C_a=\begin{pmatrix}0&0\\[2pt]0&1\end{pmatrix}
 \quad(a\in\F_q^\times).
\end{equation}

\begin{proposition}[Exact cylinder fibers]\label{prop:fibers}
For $w=w_0\cdots w_{n-1}\in\F_q^n$, the number of hidden words
$u_0\cdots u_n\in\F_q^{n+1}$ satisfying
$u_i u_{i+1}=w_i$ for $0\leq i<n$ is
\begin{equation}\label{eq:fiber-count}
 N_q(w)=(1,m)C_{w_0}C_{w_1}\cdots C_{w_{n-1}}
 \binom11.
\end{equation}
Consequently
\begin{equation}\label{eq:cylinder-law}
 \mu_q([w])=q^{-(n+1)}N_q(w).
\end{equation}
\end{proposition}

\begin{proof}
Suppose a collection of partial lifts contains $z$ lifts whose current
hidden endpoint is zero and $r$ whose endpoint is nonzero.  If the next
observed symbol is zero, a zero endpoint has one zero and $m$ nonzero
extensions, while a nonzero endpoint has only the zero extension.  Thus the
new counts are $(z+r,mz)=(z,r)C_0$.  If the next observed symbol is a fixed
$a\ne0$, no zero endpoint extends.  Each nonzero endpoint $u$ has the unique
extension $a/u$, so the new counts are $(0,r)=(z,r)C_a$.  Before any edge is
read there is one zero and $m$ nonzero choices, giving the initial row
$(1,m)$.  Summing the two terminal classes proves \eqref{eq:fiber-count}.
Uniformity of $U_0,\ldots,U_n$ gives \eqref{eq:cylinder-law}.
\end{proof}

The formula depends on a nonzero symbol only through its being nonzero.
Actual nonzero labels nevertheless remain distinct in the observed
language; their multiplicity will enter the topological transfer matrix.

\section{Exact support and word complexity}
\label{sec:support}

Let $\F_q^\times=\F_q\setminus\{0\}$.  Denote by $\cX_q$ the two-step shift
of finite type over $\F_q$ whose forbidden list is
\begin{equation}\label{eq:forbidden}
 \mathcal F_q=\{a0b:a,b\in\F_q^\times\}.
\end{equation}

\begin{theorem}[Support theorem]\label{thm:support}
For every prime power $q\geq2$,
\[
 X_q=\Phi(\F_q^\Z)=\cX_q.
\]
Every nonempty cylinder in $X_q$ has positive $\mu_q$-measure.  The shift
$X_q$ is topologically mixing.
\end{theorem}

\begin{proof}
If $y_{i-1}=a\ne0$, then $u_i\ne0$.  If also $y_i=0$, the equality
$u_i u_{i+1}=0$ forces $u_{i+1}=0$, and hence $y_{i+1}=0$.  Thus no image
contains a word $a0b$ with both exterior symbols nonzero.

Conversely, take $y\in\cX_q$.  On every maximal interval of nonzero edges,
choose one endpoint in $\F_q^\times$ and solve successively by division.
Two such intervals are separated by at least two zero edges.  The first
zero edge after a nonzero interval can therefore be lifted by making its new
endpoint zero, and the last zero edge before the next interval is then
automatically compatible.  Set any remaining unconstrained endpoints in a
zero gap equal to zero.  The same construction works for one-sided or
bi-infinite nonzero intervals; for an all-nonzero configuration choose one
hidden value and solve in both directions.  This constructs
$u\in\F_q^\Z$ with $\Phi(u)=y$.

Proposition \ref{prop:fibers} now shows that every legal finite word has a
positive cylinder probability.  Finally, two legal words can be joined by
two zeros: the only newly created length-three words meeting the bridge have
at least two consecutive zeros.  Arbitrarily many further zeros may be
inserted, so the shift is topologically mixing.
\end{proof}

For exact word counts, collapse each symbol to its zero/nonzero type.  In the
ordered pair states $00,0*,*0,**$, where $*$ means nonzero, appending a
nonzero symbol has multiplicity $m$.  The resulting weighted transfer matrix
is
\begin{equation}\label{eq:transfer}
 A_q=
 \begin{pmatrix}
 1&m&0&0\\
 0&0&1&m\\
 1&0&0&0\\
 0&0&1&m
 \end{pmatrix}.
\end{equation}
Its directed graph is strongly connected and has a loop, hence $A_q$ is
primitive.  A direct determinant calculation gives
\begin{equation}\label{eq:charpoly}
 \det(tI-A_q)
 =t\bigl(t^3-qt^2+mt-m\bigr).
\end{equation}

\begin{theorem}[Complexity recurrence and support entropy]
\label{thm:complexity}
Let $L_n=|\cL_n(X_q)|$, with $L_0=1$.  Then
\begin{align}
 L_1&=q, & L_2&=q^2, & L_3&=q^3-(q-1)^2,
 \label{eq:initial-complexity}\\
 L_{n+3}&=qL_{n+2}-(q-1)L_{n+1}+(q-1)L_n
 \quad(n\geq0).\label{eq:complexity-recurrence}
\end{align}
If $\lambda_q$ is the largest real root of
\begin{equation}\label{eq:pf-cubic}
 p_q(t)=t^3-qt^2+(q-1)t-(q-1),
\end{equation}
then
\begin{equation}\label{eq:top-entropy}
 \htop(X_q)=\log\lambda_q.
\end{equation}
\end{theorem}

\begin{proof}
Words of length at most two are unrestricted, while the forbidden list
contains $m^2$ length-three words.  This gives
\eqref{eq:initial-complexity}.  For $n\geq2$, all length-$n$ words are counted
by starting with their weighted pair state and applying $A_q^{n-2}$.
Cayley--Hamilton and \eqref{eq:charpoly} give
\eqref{eq:complexity-recurrence} for $n\geq3$.  The three remaining cases are
checked directly.  Besides $L_3=q^3-m^2$, inclusion--exclusion over the
possible locations of a forbidden word gives
\[
 L_4=q^4-2qm^2,
 \qquad
 L_5=q^5-3q^2m^2+m^3.
\]
Indeed, adjacent forbidden occurrences are incompatible, and at length five
the first and third locations overlap in exactly $m^3$ words.  Substitution
of $L_0,\ldots,L_5$ verifies the recurrence for $n=0,1,2$.
Primitivity and Perron--Frobenius give
$L_n=\exp(n\log\lambda_q+O(1))$, proving \eqref{eq:top-entropy}.
\end{proof}

For example, $L_n$ for $q=2$ begins
\[
 1,2,4,7,12,21,37,65,114,200,
\]
and for $q=3$ it begins
\[
 1,3,9,23,57,143,361.
\]
These are definition-level counts, not fitted recurrences.

\section{Short dependence and unbounded prediction memory}
\label{sec:memory}

\begin{proposition}[Stationarity, reversal, and one-dependence]
\label{prop:mixing}
The process $Y$ is stationary and time-reversible.  Moreover,
\[
 \sigma(Y_i:i\leq0)
 \quad\text{and}\quad
 \sigma(Y_i:i\geq2)
\]
are independent.  Thus $Y$ is one-dependent, strongly mixing with mixing
coefficient zero from lag two onward, and ergodic.
\end{proposition}

\begin{proof}
The code \eqref{eq:factor} commutes with the shift, so stationarity follows
from stationarity of the iid source.  Put $V_i=U_{1-i}$.  Then $V$ has the
same law as $U$, and commutativity of field multiplication gives
$V_iV_{i+1}=U_{1-i}U_{-i}=Y_{-i}$.  Hence the process is reversible.

The past sigma-field displayed in the statement is contained in
$\sigma(U_i:i\leq1)$, while the future sigma-field is contained in
$\sigma(U_i:i\geq2)$.  These source sigma-fields are independent.  The
mixing and ergodicity conclusions follow.
\end{proof}

Short dependence does not bound the length of a predictive context.  Define
the generalized Fibonacci sequence
\begin{equation}\label{eq:F-recurrence}
 F_0=0,\qquad F_1=1,\qquad
 F_{r+2}=F_{r+1}+mF_r\quad(r\geq0).
\end{equation}
For $a\in\F_q^\times$ and $r\geq1$, let
\begin{equation}\label{eq:context-event}
 E_r(a)=\{Y_{-r-1}=a,\ Y_{-r}=\cdots=Y_{-1}=0\}.
\end{equation}
Each event has positive probability.

\begin{theorem}[Exact contexts and infinite Markov order]
\label{thm:infinite-memory}
For $a,b\in\F_q^\times$ and $r\geq1$,
\begin{equation}\label{eq:context-probability}
 \Prob(Y_0=b\mid E_r(a))
 =\frac{mF_{r-1}}{qF_{r+1}}.
\end{equation}
The value is independent of $a$ and $b$.  As $r$ increases, the right-hand
side never takes the same value at two consecutive indices and alternates
around its limit.  Consequently $\mu_q$ is not a Markov measure of any
finite order.
\end{theorem}

\begin{proof}
Conditioning on $Y_{-r-1}=a\ne0$ makes the right hidden endpoint of that edge
nonzero.  Starting from one nonzero endpoint, the zero-edge transition in
\eqref{eq:fiber-matrices} gives
\begin{equation}\label{eq:zero-run-counts}
 (0,1)C_0^r=(F_r,mF_{r-1}).
\end{equation}
The identity follows by induction from \eqref{eq:F-recurrence}.  There are
therefore $F_{r+1}=F_r+mF_{r-1}$ compatible continuations through the zero
run, of which $mF_{r-1}$ end at a nonzero value $U_0$.  Conditional on such
a nonzero $U_0$, exactly one of the $q$ equally likely choices for $U_1$
satisfies $U_0U_1=b$.  This proves \eqref{eq:context-probability}.

The same calculation is unchanged by an arbitrary compatible earlier word
$v$.  Indeed, if $(1,m)C_v=(z,c)$, then reading $a\ne0$ gives
$(z,c)C_a=(0,c)$; the scalar $c$ multiplies both the numerator and denominator
after the zero run and cancels.  Thus \eqref{eq:context-probability} is a
genuine context conditional, not merely a calculation under a truncated past.

Set $t_r=mF_{r-1}/F_r$ for $r\geq1$.  Then $t_1=0$ and
\begin{equation}\label{eq:mobius}
 t_{r+1}=\frac{m}{1+t_r}.
\end{equation}
The map on the right is injective and decreasing on $[0,\infty)$.  Its
unique nonnegative fixed point is
\begin{equation}\label{eq:tstar}
 t_*=\frac{\sqrt{4q-3}-1}{2}.
\end{equation}
No finite iterate of $t_1=0$ can equal $t_*$: injectivity would otherwise
allow repeated backward cancellation and force $0=t_*$.  Thus
$t_{r+1}\ne t_r$ for every $r$, and the decreasing recurrence alternates
strictly around $t_*$.  Since the probability in
\eqref{eq:context-probability} equals
$t_r/[q(1+t_r)]$, the same conclusions hold for the prediction values.

If the process had Markov order $k\geq1$, all earlier histories ending in
the same $k$ zeros would give the same next-symbol law.  Taking consecutive
$r,r+1\geq k$ in \eqref{eq:context-event} contradicts this.  Order zero is
also impossible: after the allowed context $a0$ a nonzero next symbol has
probability zero, whereas its unconditional probability is positive.
\end{proof}

The first few probabilities for a fixed nonzero $b$ are
\[
 0,\quad \frac{q-1}{q^2},\quad
 \frac{q-1}{q(2q-1)},\quad
 \frac{q-1}{q^2+q-1},\quad\ldots,
\]
where the unsimplified pattern is governed by
\eqref{eq:context-probability}.  For $q=2$ the sequence begins
\[
 0,\quad \frac14,\quad \frac16,\quad \frac15,\quad
 \frac3{16},\quad \frac5{26}.
\]

\section{An exact entropy-rate series}
\label{sec:entropy}

The nonzero symbols are regeneration markers for prediction.  Let $R$ be
the number of consecutive zeros immediately before time zero:
\begin{equation}\label{eq:age}
 R=\min\{r\geq0:Y_{-r-1}\ne0\}.
\end{equation}
Since $\Prob(Y_i\ne0)=m^2/q^2>0$ and the process is ergodic, $R$ is finite
almost surely.

Define
\begin{equation}\label{eq:alpha-s}
 \alpha_0=1,
 \qquad
 \alpha_r=\frac{mF_{r-1}}{F_{r+1}}\ (r\geq1),
 \qquad
 s_r=\frac{m}{q}\alpha_r.
\end{equation}
Here $\alpha_r$ is the posterior probability that $U_0\ne0$ given a past
whose last nonzero observation is followed by $r$ zeros, and $s_r$ is the
conditional probability that $Y_0\ne0$.

Write
\[
 \hb(s)=-s\log s-(1-s)\log(1-s),
\]
with $0\log0=0$.

\begin{theorem}[Entropy series and strict entropy gap]
\label{thm:entropy}
The age distribution is
\begin{equation}\label{eq:age-law}
 w_r:=\Prob(R=r)=\frac{m^2F_{r+1}}{q^{r+2}},
 \qquad r\geq0.
\end{equation}
The entropy rate of $\mu_q$, in nats, is the exponentially convergent series
\begin{equation}\label{eq:entropy-series}
 h(\mu_q)
 =\sum_{r=0}^{\infty}w_r
 \bigl(\hb(s_r)+s_r\log m\bigr).
\end{equation}
Moreover,
\begin{equation}\label{eq:strict-gap}
 h(\mu_q)<\htop(X_q)=\log\lambda_q.
\end{equation}
\end{theorem}

\begin{proof}
The event $Y_{-r-1}\ne0$ has probability $m^2/q^2$.  Given that event,
the right endpoint of the observed edge is nonzero.  There are
\[
 (0,1)C_0^r\binom11=F_{r+1}
\]
hidden continuations through $r$ zero edges, out of $q^r$ total
continuations.  This proves \eqref{eq:age-law}.  The generating function
\begin{equation}\label{eq:F-generating}
 \sum_{j\geq0}F_jz^j=\frac{z}{1-z-mz^2}
\end{equation}
shows that $\sum_{r\geq0}w_r=1$ after setting $z=1/q$.

Conditioning further on any compatible observations before the last nonzero
edge can bias its hidden endpoint among the nonzero field elements, but it
cannot change that endpoint's zero/nonzero type.  Every nonzero starting value
has the same type-counts through a zero run, so the common scalar cancels from
the posterior ratio.  Consequently, given the complete observed past with
$R=r$, the calculation in \eqref{eq:zero-run-counts} gives
$\Prob(U_0\ne0\mid Y_{-\infty}^{-1})=\alpha_r$.  If $U_0\ne0$, then the
independent uniform $U_1$ makes each of the $m$ nonzero values of $Y_0$
have probability $1/q$.  Hence, conditional on $R=r$, the symbol zero has
probability $1-s_r$ and each nonzero symbol has probability $s_r/m$.
Its entropy is $\hb(s_r)+s_r\log m$.  Averaging over $R$ proves
\eqref{eq:entropy-series} by the standard identity
$h(\mu_q)=H(Y_0\mid Y_{-\infty}^{-1})$.

The growth rate of $F_r$ is the positive root of $x^2=x+m$, which is
strictly smaller than $q$; indeed $q^2-q-m=m^2>0$.  Thus the weights in
\eqref{eq:age-law} decay exponentially.

By \cref{thm:support}, $X_q$ is a mixing shift of finite type.  Its unique
measure of maximal entropy is the finite-state intrinsic Markov measure
constructed by Parry \cite{Parry1964}.  On the ordered-pair presentation its
transition law is one-step Markov, so in the original symbol coordinates it
has Markov order at most two.  Theorem \ref{thm:infinite-memory} shows that
$\mu_q$ is not that measure.  Uniqueness of the maximal measure therefore
gives the strict inequality in \eqref{eq:strict-gap}.
\end{proof}

For orientation, the control program evaluates
\begin{center}
\begin{tabular}{c@{\qquad}cc}
\toprule
$q$ & $h(\mu_q)$ (nats) & $\log\lambda_q$ (nats)\\
\midrule
2 & 0.484678454953871 & 0.562399148645924\\
3 & 0.853898202937607 & 0.924806254398638\\
4 & 1.134074691264723 & 1.216224572920358\\
5 & 1.361005489603169 & 1.454951935105873\\
\bottomrule
\end{tabular}
\end{center}
The strict gap is structural; the decimals are not used in its proof.

\section{Ownership boundary and scope}
\label{sec:ownership}

Aaronson, Gilat, Keane, and de Valk \cite{AaronsonEtAl1989} belong to the
general one-dependent/block-factor ownership line.  Blackwell's work
\cite{Blackwell1957} owns the general problem of entropy for functions of
finite-state Markov chains.  Rissanen
\cite{Rissanen1983} introduced context-based variable-memory models, and
B\"uhlmann and Wyner \cite{BuhlmannWyner1999} developed variable-length
Markov chains as a statistical class.  Parry \cite{Parry1964} owns the
intrinsic Markov measure used only to identify the strict entropy gap.

The present note does not claim any of these frameworks.  Its theorem
package is the calculation for the specific singular edge factor
$u_i u_{i+1}$ over $\F_q$: the forbidden support, exact fiber matrices,
complexity cubic, generalized-Fibonacci context probabilities, and entropy
series.  A bounded search through 28 August 2026 using the exact defining
formula and combinations of ``adjacent product'', ``finite field'',
``one-dependent'', and ``hidden Markov'' found no direct primary source for
this combined package.  That negative search is a narrow collision firewall,
not evidence of absolute priority.

Several extensions are deliberately left open.  Replacing the field by a
finite ring introduces zero divisors and more than two endpoint types;
replacing multiplication by a general binary operation asks which
singularities create an unbounded context tree.  Both directions require new
fiber semigroups and are not consequences of the two matrices above.

\section{Conclusion}

Adjacent products over a finite field give a compact example in which local
dependence ends after one site but predictive memory is unbounded.  The
field zero turns the image into a mixing two-step SFT, while uncertainty
inside zero runs produces a generalized-Fibonacci context law.  The same
two-state fiber calculation controls cylinders, word growth, failure of
finite Markov order, and the entropy rate.  The example isolates a reusable
mechanism: a finite-block factor can be strongly mixing at a fixed finite
range and still retain arbitrarily long variable-length memory through a
singular observation symbol.

\bibliographystyle{plainnat}
\bibliography{references}

\end{document}
